% Generated draft: requires mathematical and source review.
\documentclass[11pt]{article}
\usepackage[margin=25mm]{geometry}
\usepackage{fontspec}
\setmainfont{Latin Modern Roman}
\usepackage{amsmath,amssymb,mathtools,unicode-math}
\setmathfont{Latin Modern Math}
\usepackage{enumitem,needspace,xurl,hyperref}
\setlength{\parindent}{0pt}
\setlength{\parskip}{5pt}
\setlength{\emergencystretch}{3em}
\newcommand{\N}{\mathbb{N}}
\newcommand{\Z}{\mathbb{Z}}
\newcommand{\Q}{\mathbb{Q}}
\newcommand{\R}{\mathbb{R}}
\newcommand{\C}{\mathbb{C}}
\newcommand{\F}{\mathbb{F}}
\newcommand{\A}{\mathbb{A}}
\newcommand{\PP}{\mathbb{P}}
\newcommand{\E}{\mathbb{E}}
\newcommand{\eps}{\varepsilon}
\newcommand{\dd}{\,\mathrm{d}}
\begin{document}
\section*{UnsolvedMath ID: 30004118}
\textbf{Source problem code:} catalogue entry
\subsection{Definitions and mathematical statement}
Let \(\mathbb H^n\) denote real hyperbolic \(n\)-space of sectional curvature \(-1\), and let \(\operatorname{PO}(n,1)\) be its projective orthogonal isometry group. A lattice in this group is a discrete subgroup \(\Gamma\) for which \(\Gamma\backslash\operatorname{PO}(n,1)\) has finite Haar volume. A real hyperbolic \(n\)-manifold is a Riemannian \(n\)-manifold of constant sectional curvature \(-1\). Completeness and finite volume are additional conditions that the supplied manifold question does not explicitly state. A submanifold is totally geodesic if its second fundamental form vanishes. A number field \(k\) is totally real if every field embedding \(\sigma:k\hookrightarrow\C\) has image in \(\R\).

The corrected source record asks the following two questions.
\begin{enumerate}
\item Does there exist an integer \(n\geq4\) and a real hyperbolic \(n\)-manifold \(M\) such that, for every integer \(m\) satisfying \(2\leq m\leq n-1\), \(M\) contains no arithmetic totally geodesic \(m\)-submanifold?

Here ``arithmetic'' is understood as referring to arithmeticity of the hyperbolic lattice defining the submanifold, rather than merely to arithmeticity of an abstract group. This interpretation requires clarification of the intended submanifolds: the record does not specify whether they must be complete, have finite volume, or be embedded rather than immersed. It also does not specify completeness or finite volume of \(M\). Thus these conditions must be fixed before this question has a fully determined formulation in terms of hyperbolic lattices. The displayed existential quantifier follows the wording ``a hyperbolic \(n\)-manifold, with \(n\geq4\)''; the record does not explicitly request an example for each such dimension.
\item For every integer \(n\geq4\) and every lattice \(\Gamma\leq\operatorname{PO}(n,1)\), is the trace field of \(\Gamma\) totally real?

The phrase ``trace field'' is not defined in the supplied record. Two possible precise conventions are the field of traces in the defining orthogonal representation and the adjoint trace field. To describe these without choosing between them, let \(\pi:\operatorname{O}(n,1)\to\operatorname{PO}(n,1)\) be the projectivization map and set
\[
 k_{\mathrm{def}}(\Gamma)=\Q\bigl(\operatorname{tr}(A):A\in\pi^{-1}(\Gamma)\bigr),
 \qquad
 k_{\mathrm{Ad}}(\Gamma)=\Q\bigl(\operatorname{tr}(\operatorname{Ad}(\gamma)):\gamma\in\Gamma\bigr),
\]
where \(\operatorname{Ad}\) acts on the real Lie algebra \(\mathfrak{so}(n,1)\). Using the full inverse image in the first expression avoids choosing representatives for projective elements. The intended trace convention must be confirmed; no equivalence of the resulting questions is asserted here. The source also leaves implicit the number-field setting needed for the stated definition of total reality.
\end{enumerate}

The original statement additionally asked whether there are infinitely many commensurability classes of nonarithmetic lattices in \(\operatorname{PU}(2,1)\), the projective unitary isometry group of complex hyperbolic two-space. The supplied correction explicitly removes that clause because it is subsumed by source problem \texttt{OWR-16929-001}. It is retained here as provenance, not as a third target for this entry. The supplied open and verification labels are registry metadata, not independent mathematical verification.
\subsection{Short English statement}
The corrected entry asks whether some real hyperbolic manifold of dimension at least four has no arithmetic totally geodesic submanifold of intermediate dimension, and whether every lattice in \(\operatorname{PO}(n,1)\), for every \(n\geq4\), has totally real trace field. The intended completeness, volume, embedding, and trace-field conventions require clarification. A former question about \(\operatorname{PU}(2,1)\) was removed as a duplicate.
\subsection{Sources}
\begin{itemize}
\item Supplied local UnsolvedMath registry record, permanent ID \(30004118\), source problem code \texttt{catalogue entry}, entitled ``Arithmeticity Problems for Hyperbolic Lattices and Submanifolds,'' in the Oberwolfach Reports Open Problems collection. The record lists proposed year 2019 and no named proposer. Its supplied status is \texttt{open}; its statement-status label is \texttt{corrected\_verified}. The accompanying verification note says the \(\operatorname{PU}(2,1)\) clause was removed because it is subsumed by \texttt{OWR-16929-001}.
\item Supplied source citation: \emph{Mini-Workshop: Reflection Groups in Negative Curvature} (2020), Oberwolfach Reports, pp. 1043--1070, DOI 10.4171/owr/2019/17. \url{https://doi.org/10.4171/owr/2019/17}. The supplied background describes Oberwolfach Reports as documenting mathematical workshops and their open-problem discussions.
\item Supplied literature source: ``Oberwolfach Report 19/2019, selected open problems.'' \url{https://ems.press/content/serial-article-files/46796}.
\item Supplied literature source: Deraux, \emph{A new nonarithmetic lattice in PU(3,1)}. \url{https://doi.org/10.2140/agt.2020.20.925}.
\item Literature assessment dated exactly \texttt{2026-08-22}: the supplied curation labels the record open and states that three questions about nonarithmetic complex or real hyperbolic lattices remain open in the workshop report and are not settled by later constructions. This is supplied dated curation, not independently verified. Its reference to three questions reflects the original wording, whereas the corrected statement retains two. The record itself requests independent verification of time-sensitive status and cited claims.
\end{itemize}
\end{document}
